\section{WC 11/11/24}

\paragraph{Overview:}
This week I took my project details and started a technical dive into the various skills I would need to acquire to make my project feasible, I also continued to think about the problems I wanted to solve with my artifact

\paragraph{Achievements, Thoughts and Ideas:}
This week I further explored my current research and searched for some new content to cover. My project has already spread in so many interesting directions that it was crucial to ensure my sources were well-organised, documented and valid.

Additionally, I have added some podcasts to my reading list that I can work my way through during times I wouldn't usually be able to complete work (such as when running), by listening to audio and making brief mental notes I can then sit down to formally document my learnings with greater insight.

I have also looked into conducting some first party research in 2 separate forms. Gaining wider public feedback once I have a prototype version of my application - to help direct design decisions, and to interview some leading researchers and scientists in the industry of memory techniques and human data comprehension, I have already started some discussion to make these a reality and am planning an interview with a leading researcher and author in the field of self-improvement and learning

I continued to dive into the huge variety of technologies and skills required to create a web application so that I can understand some technical limitations which may impact my design considerations, some examples of this include: \nameref{research:DataTagging} - a technique that can help link the human technique of mind mapping to a computer system, \nameref{research:WebBasedDrawing} - a feature I feel is key to the success of my project, \nameref{research:DatabaseStorageMethods} - it would be pointless creating vast arrays of notes if they could be stored and \nameref{research:LiveUserCollaboration} - core techniques that allow my consumer facing webapps to interface with my backend data storage which allows for multi user collaboration and data sharing

Whilst researching these tools I started building a couple of small prototypes (\nameref{prototype:LiveUserCollaboration}, \nameref{prototype:WebDrawing}), drawing both from my work experience learnings and my prior research. Although these are simple tools, they show that the individual key elements of my project are viable, and I have built a more advanced tool that will hopefully be functioning midway next week

I also learnt about a platform called Figma which is a user interface design tool, after following a course I created some \nameref{research:FigmaDesign} projects to become familiar with Figma's feature set before continuing further into my own applications design journey

\paragraph{What went well?}
I now have a very thorough understanding of the individual components required to produce my end software and now just need too apply these learnings before finding other areas to research


\paragraph{What was challenging?}
After a quick conversation with my computer science teacher he raised some potential concerns with my Software based artifact overlapping the computer science NEA coursework which would follow in year 13.
\paragraph{What was done to resolve this?}
I brought my concerns to my EP mentor who sat down with me to discuss the similarity and subject crossover. I explained that all of the technologies I would be using in this project are not mentioned / taught in the course. We also discussed my goals of using the software artifact to quantify and realise my learnings of human learning methods not just using the project to learn about computational techniques. With the help of my mentor we concluded that this project would be viable but that I would need to be clear in my separation of the project from computer science NEA by thoroughly recording all techniques learnt and all research regarding human learning
\paragraph{Plans for next week?}
Next week I hope to have my note exploration tool finalised in addition to cleaning up parts of my diary and research notes. Until this point I have been focusing on the content of my learnings but writing up my thoughts and takeaways is just as important

